\subsection{Single-Episode Trajectory Analysis}

Figure~\ref{fig:single_episode} shows a representative evaluation episode for the trained
DQN agent on the 20-sensor, 500$\times$500 condition. The agent achieves 100\% sensor
coverage by step~200 and sustains data collection until battery depletion at step~1{,}404,
collecting 57{,}335\,bytes at a mean efficiency of 213.6\,B/Wh. Cumulative reward grows
monotonically without plateauing, indicating that the agent does not idle once early sensors
are visited---a direct consequence of the buffer-overflow penalty encouraging continued
collection from revisited sensors.

\begin{figure}[H]
    \centering
    \includegraphics[width=0.68\textwidth]{sensors_20/fig1_shaded_reward}
    \caption{Cumulative reward and battery level during a single DQN evaluation episode
             (20 sensors, 500$\times$500 grid, seed 42). Shaded band shows $\pm 1$ standard
             deviation across seeds.}
    \label{fig:single_episode}
\end{figure}

\subsubsection{Spatial Trajectory Comparison}

Figure~\ref{fig:trajectory_comparison} shows the UAV flight paths for all three agents on
the same sensor layout (seed 42, $N=20$), with sensor positions coloured by their initial SF
assignment. The contrast is qualitatively striking.

The \textbf{Nearest Greedy} agent executes a near-straight sweep from sensor to sensor
in order of proximity, issuing a collection action at every step (100\% hover fraction).
The \textbf{SF-Aware Greedy} agent follows a similar sweeping path but visits higher-priority
(SF7--SF9) sensors first; like Nearest Greedy, it never repositions without an immediate
collection target.

The \textbf{DQN} agent, by contrast, devotes a measurable fraction of steps (2.9\% across
10 seeds, $p = 0.002$, Welch's $t$-test) to deliberate movement---visible as the orange hover
dots dispersed across the grid rather than concentrated on sensor locations. The DQN's
trajectory is less linear: it navigates towards spatial positions that simultaneously
place multiple sensors within collection range across \emph{orthogonal} spreading factors,
rather than visiting each sensor independently.

\begin{figure}[H]
    \centering
    \includegraphics[width=0.82\textwidth]{fig_trajectory_comparison}
    \caption{UAV trajectories for all three agents on the same sensor layout (seed 42,
             $N=20$, 500$\times$500 grid). Sensor markers are coloured by initial SF
             assignment (green = SF7, purple = SF12). Orange dots = hover steps; coloured
             dots = movement steps. The DQN's deliberate repositioning (orange cluster
             spread across the grid) contrasts with the greedy agents' end-to-end sweeps.}
    \label{fig:trajectory_comparison}
\end{figure}

\subsubsection{Spreading Factor Distribution Over Time}

Figure~\ref{fig:sf_dynamics} shows how the distribution of spreading factors across the
20 sensors evolves throughout the episode for the DQN and SF-Aware Greedy agents. Two
observations are notable.

The DQN systematically reduces the proportion of sensors on high SFs (SF11--SF12) by
maintaining proximity for multiple steps, allowing the EMA-ADR filter to converge to lower
SF assignments before repositioning. This maintains greater SF diversity (SF7--SF11) for
longer than the SF-Aware Greedy baseline, confirming that deliberate repositioning actively
manages the spectral environment rather than wasting movement budget. The causal mechanism
is analysed in detail in Section~\ref{sec:discussion}.

\begin{figure}[H]
    \centering
    \includegraphics[width=0.72\textwidth]{fig_sf_dynamics}
    \caption{Spreading factor distribution across all 20 sensors over the episode (seed 42).
             The DQN (left) maintains greater SF diversity and achieves faster SF downgrades
             than the SF-Aware Greedy baseline (right), confirming that deliberate
             repositioning actively reshapes the LoRaWAN spectral environment.}
    \label{fig:sf_dynamics}
\end{figure}
